\section{Contexto y datos}

\subsection{Contexto del evento}

Se considera un evento academico de ciencia de datos organizado en dos dias, con ocho franjas temporales y cuatro franjas por dia. La organizacion debe decidir que actividades ofrece, cada una como mucho una vez, y asignar su inicio a un espacio compatible y disponible. Las decisiones afectan a la asistencia esperada, a la valoracion de las actividades, al coste total y al numero de actividades ofrecidas.

Las actividades son conferencias, talleres y mesas redondas. Las franjas representan bloques horarios discretos. Una actividad de duracion $d_i$ ocupa las franjas $t,t+1,\ldots,t+d_i-1$. No puede superar la franja 8 ni cambiar de dia.

\subsection{Instancia principal}

La instancia contiene 18 actividades, tres espacios, ocho franjas y tres categorias. Hay dos actividades obligatorias, talleres de varias franjas y relaciones de incompatibilidad. Los costes se expresan en unidades monetarias y el presupuesto principal es 6500.

\begin{table}[h]
\centering
\begin{tabular}{llrr}
\hline
Espacio & Nombre & Capacidad & Franjas disponibles \\
\hline
E1 & Auditorio & 200 & 1--8 \\
E2 & Aula de ordenadores & 40 & 1,2,3,5,6,7,8 \\
E3 & Sala de seminarios & 80 & 2,3,4,6,7,8 \\
\hline
\end{tabular}
\caption{Espacios de la instancia principal.}
\end{table}

Los datos estan separados del codigo en `data/`. Las relaciones se mantienen mediante identificadores. Una fila sin disponibilidad significa que no existe esa franja disponible para la actividad o el espacio correspondiente.

\subsection{Conjuntos y parametros}

\begin{itemize}
\item $I$: actividades.
\item $E$: espacios.
\item $T=\{1,\ldots,8\}$: franjas.
\item $d_i$: duracion de $i$ en franjas.
\item $a_i$: asistencia esperada.
\item $v_i$: valoracion.
\item $c_i$: coste.
\item $B$: presupuesto.
\item $\mathrm{cap}_e$: capacidad de $e$.
\item $F=8$: horizonte; $FD=4$: franjas por dia.
\item $w_A,w_V,w_N$: pesos normalizados de asistencia, valoracion y numero de actividades.
\end{itemize}

\subsection{Variables}

Se usa una variable binaria
\[
y_i=
\begin{cases}
1,&\text{si se programa la actividad }i,\\
0,&\text{en otro caso.}
\end{cases}
\]

Se usa una variable binaria $x_{i,e,t}$ que vale 1 si la actividad $i$ comienza en el espacio $e$ en la franja $t$.

\subsection{Restricciones}

\paragraph{Seleccion unica.}
\[
\sum_{e\in E}\sum_{t\in T}x_{i,e,t}=y_i
\qquad\forall i\in I.
\]

\paragraph{Actividades obligatorias.}
\[
y_i=1\qquad\forall i\in O.
\]

\paragraph{Compatibilidad y capacidad.}
Solo se crean variables $x_{i,e,t}$ para pares compatibles y con $\mathrm{cap}_e\geq a_i$.

\paragraph{Disponibilidad y horizonte.}
Solo se crean variables cuya actividad y espacio estan disponibles en todas las franjas ocupadas. Tambien debe cumplirse $t+d_i-1\leq F$ y que $t$ y $t+d_i-1$ pertenezcan al mismo dia.

\paragraph{Presupuesto.}
\[
\sum_{i\in I}c_i y_i\leq B.
\]

\paragraph{Ocupacion de espacios.}
Para cada $e$ y franja $s$,
\[
\sum_{i\in I}\sum_{\substack{t\in T\\t\leq s<t+d_i}}
x_{i,e,t}\leq 1.
\]
Esta restriccion incluye en el solapamiento las franjas interiores de una actividad de varias franjas.

\subsection{Funcion objetivo}

Se maximiza
\[
Z=
w_A\frac{\sum_i a_i y_i}{\sum_i a_i}
+
w_V\frac{\sum_i v_i y_i}{\sum_i v_i}
+
w_N\frac{\sum_i y_i}{|I|},
\]
con $w_A+w_V+w_N=1$. Las tres componentes estan acotadas entre cero y uno. Esta normalizacion evita que un criterio domine solo por la escala numerica. Para priorizar un criterio se fija su peso en uno y los otros en cero.

\subsection{Ampliacion}

Los pares $(i,j)\in P$ son incompatibles y no pueden solaparse. Para cada $s\in T$ se impone
\[
\sum_{\substack{t\in T\\t\leq s<t+d_i}}\sum_{e\in E}x_{i,e,t}
+
\sum_{\substack{t\in T\\t\leq s<t+d_j}}\sum_{e\in E}x_{j,e,t}
\leq 1.
\]
